\chapter{Introduction}
\label{chap:introduction}

The increasing availability of high-throughput genotyping
technologies has made it possible to characterize livestock
populations using tens of thousands of Single Nucleotide
Polymorphisms (SNPs)
\cite{georges2019livestock, morin2004snps}.

From a computational perspective, these datasets pose a distinctive
high-dimensional learning problem.
The number of genomic variables can substantially exceed the number
of available individuals, producing the well-known $p \gg n$
setting
\cite{saeys2007feature, chafai2023mlanimal}.

This imbalance is particularly relevant in livestock genomics,
where some breeds can be represented by relatively few genotyped
individuals while each individual is described by tens of thousands
of markers
\cite{groeneveld2010diversity, chafai2023mlanimal}.
At the same time, genomic features are not statistically
independent, since linkage disequilibrium can introduce substantial
correlation between nearby loci
\cite{slatkin2008linkage, pritchard2001linkage}.

These characteristics raise two practical computational questions.
First, how many individuals per breed are required to train a
reliable classification model?
Second, how many SNP markers are actually required to retain most
of the breed-discriminative information available in the original
high-density genomic representation?

This thesis investigates these questions using classical machine
learning, deep representation learning, and Explainable Artificial
Intelligence (XAI).
Rather than assuming a predefined minimum sample size or SNP-panel
size, both quantities are treated as experimental variables and are
systematically evaluated in Chapter~\ref{chap:experiments}.

A further objective is to determine whether marker rankings derived
from deep models provide useful information for genomic
dimensionality reduction when compared with established
population-genetics and machine-learning approaches.
The feature-selection methodology is described in
Chapter~\ref{chap:feature_importance}.


\section{Inductive Bias in Computational Genomics}
\label{sec:inductive_bias_genomics}

Learning algorithms rely on explicit or implicit assumptions about
the structure of the data in order to generalize beyond the
observations available during training.
These assumptions are commonly described as the
\emph{inductive bias} of a learning algorithm
\cite{mitchell1980need}.

Different model families therefore emphasize different properties
of genomic data.
Classical locus-wise population-genetics statistics, for example,
evaluate individual markers according to quantities such as allele
frequencies or between-population differentiation
\cite{wright1949genetical, weir1984estimating}.

Genome-Wide Association Studies (GWAS) traditionally evaluate
associations between individual genomic variants and a phenotype
or trait of interest
\cite{hirschhorn2005genome}.

The locus-wise formulation can be challenging when the number of
available individuals is small relative to the number of genomic
markers. In cattle, Manca et al.\ investigated a multivariate GWAS
framework based on discriminant analysis and explicitly evaluated
its behavior across increasing sample sizes
\cite{manca2020multivariate}.
Although phenotype association is distinct from the breed-assignment
task addressed in this thesis, that work provides a relevant
methodological precedent for treating both sample size and joint
marker structure as important experimental factors.

Multivariate approaches instead consider multiple markers jointly.
In livestock genomics, discriminant methods have been used for
breed assignment, reduced-marker selection, and the analysis of
population differentiation in both cattle and sheep
\cite{dimauro2013canonical, dimauro2015sheep,
      sorbolini2016signatures}.


Artificial Neural Networks introduce an additional non-linear
inductive bias through compositions of learned transformations.
This increased expressive capacity may be useful for modeling
complex structure in genomic datasets, but it does not by itself
guarantee superior classification or feature-selection performance
\cite{chafai2023mlanimal}.

Whether deep representations provide an advantage over classical
approaches is therefore treated as an empirical question in this
thesis and is evaluated using a common downstream protocol in
Chapter~\ref{chap:experiments}.

The conceptual distinction between locus-wise genomic analysis and
multivariate representation learning is illustrated in
Figure~\ref{fig:inductive_bias_comparison}.

\begin{figure}[htbp]
\centering
\begin{tikzpicture}[
    >=Latex,
    font=\small,
    box/.style={
        draw,
        semithick,
        rounded corners=2pt,
        align=center,
        minimum height=10mm,
        minimum width=39mm,
        inner sep=4pt,
        fill=black!2
    },
    key/.style={
        box,
        fill=black!7
    },
    marker/.style={
        draw,
        semithick,
        rounded corners=2pt,
        align=center,
        minimum height=9mm,
        minimum width=22mm,
        inner sep=3pt,
        fill=black!2
    },
    arr/.style={->, semithick},
    note/.style={
        align=center,
        font=\scriptsize,
        text=black!70
    }
]

% ------------------------------------------------------------------
% Titles
% ------------------------------------------------------------------

\node[font=\bfseries\normalsize] at (-4.1,3.45)
    {Locus-wise analysis};

\node[font=\bfseries\normalsize] at (4.1,3.45)
    {Multivariate modeling};

% ------------------------------------------------------------------
% LEFT: LOCUS-WISE
% ------------------------------------------------------------------

\node[key] (lm) at (-4.1,2.45) {
    Genotype matrix\\
    $N$ individuals $\times$ $P$ SNP markers
};

\node[note] at (-4.1,1.80) {
    each SNP is considered separately
};

\node[marker] (s1) at (-6.1,0.85) {SNP$_1$};
\node[marker] (s2) at (-4.1,0.85) {SNP$_2$};
\node[marker] (sp) at (-2.1,0.85) {SNP$_P$};

\node[box] (lst) at (-4.1,-0.70) {
    Marker-specific statistic\\
    e.g.\ $F_{ST}$ or association score
};

\node[key] (lr) at (-4.1,-2.20) {
    SNP ranking\\
    $s_1,\ldots,s_P$
};

\draw[arr] (lm) -- (s1);
\draw[arr] (lm) -- (s2);
\draw[arr] (lm) -- (sp);

\draw[arr] (s1) -- (lst);
\draw[arr] (s2) -- (lst);
\draw[arr] (sp) -- (lst);

\draw[arr] (lst) -- (lr);

% ------------------------------------------------------------------
% RIGHT: MULTIVARIATE
% ------------------------------------------------------------------

\node[key] (rv) at (4.1,2.45) {
    Individual genotype profile\\
    $\mathbf{x}\in\mathbb{R}^{P}$\\
    \scriptsize $P$ SNP values
};

\node[note] at (4.1,1.60) {
    all SNPs are processed jointly
};

\node[box] (rm) at (4.1,0.60) {
    Multivariate / neural model\\
    $f(\mathbf{x})$
};

\node[box] (rz) at (4.1,-0.85) {
    Learned representation\\
    $\mathbf{z}\in\mathbb{R}^{d}$,\quad $d \ll P$
};

\node[key] (ro) at (4.1,-2.35) {
    Breed classification\\
    and/or SNP attribution
};

\draw[arr] (rv) -- (rm);
\draw[arr] (rm) -- (rz);
\draw[arr] (rz) -- (ro);

% ------------------------------------------------------------------
% Separator
% ------------------------------------------------------------------

\draw[gray, dashed] (0,3.10) -- (0,-2.95);

% ------------------------------------------------------------------
% Bottom notation
% ------------------------------------------------------------------

\node[note] at (0,-3.35) {
    $N$: number of individuals \qquad
    $P$: number of SNP markers \qquad
    $d$: latent-space dimension
};

\end{tikzpicture}

\caption{
Conceptual comparison between locus-wise genomic analysis and
multivariate modeling.
Locus-wise approaches evaluate each of the $P$ SNP markers
individually and derive marker-specific statistics, whereas
multivariate models process the complete genotype profile jointly
and may learn a lower-dimensional representation with $d \ll P$.
$N$ denotes the number of individuals in the dataset.
Author's elaboration based on
\cite{weir1984estimating,dimauro2013canonical,bengio2013representation}.
}
\label{fig:inductive_bias_comparison}
\end{figure}




\section{Generative and Contrastive Representation Learning}
\label{sec:representation_learning}

Representation learning provides a way to replace direct reasoning in
the original high-dimensional SNP space with learned representations
that can concentrate information relevant to a downstream task
\cite{bengio2013representation}.
This perspective is particularly attractive in genomic settings,
where the input contains many correlated markers and the number of
available individuals may be comparatively small.

Two complementary paradigms are investigated in this thesis.
The first is variational representation learning, in which a
probabilistic latent variable is learned jointly with genotype
reconstruction and breed prediction
\cite{kingma2013auto, rezende2014stochastic}.
The implemented model is inspired by the Variational Multi-view
Genomics Predictor (VMGP)
\cite{zhao2025vmgp}.
Because VMGP was introduced for genomic selection in plants, its use
here is treated as an architectural adaptation to livestock breed
classification rather than as a direct transfer of an already
validated livestock model.

The second paradigm is contrastive representation learning.
Instead of reconstructing the input, contrastive objectives organize
the latent space by controlling similarity relationships between
observations or augmented views
\cite{oord2018representation, chen2020simple}.
The model considered in this thesis uses normalized latent
representations and a similarity-based objective, allowing the learned
space to be interpreted geometrically on a hypersphere
\cite{wang2020understanding}.

These paradigms impose different learning objectives and are therefore
compared empirically rather than assumed to be interchangeable.
Their theoretical foundations are reviewed in
Section~\ref{sec:deep_learning_fundamentals}, while the implemented
architectures are specified in Chapter~\ref{chap:methods}.

\section{The Black Box Problem and Explainable AI}
\label{sec:black_box_xai}

Deep neural networks can learn flexible non-linear representations,
but their internal decision processes are not directly interpretable
\cite{guidotti2018survey, arrieta2020explainable}.
In this thesis, classification performance is therefore only one part
of the analysis: a second objective is to determine which genomic
markers contribute most strongly to breed discrimination.

Post-hoc Explainable Artificial Intelligence (XAI) provides a
mechanism for assigning relevance scores to input features.
Two complementary attribution paradigms are considered: SHAP, which
is based on Shapley-value principles
\cite{shapley1953value, lundberg2017unified},
and Integrated Gradients, which exploits the differentiable structure
of neural models
\cite{sundararajan2017axiomatic}.

The sample-level attributions produced by these methods are aggregated
into global SNP rankings and evaluated through reduced marker panels.
The exact aggregation and panel-construction procedures are described
in Chapter~\ref{chap:feature_importance}.

\section{Biological and Practical Motivation}
\label{sec:biological_economic_impact}

Accurate genomic characterization of livestock populations has
applications in breed assignment, biodiversity conservation,
population management, and authentication of products of animal
origin
\cite{dimauro2013canonical, groeneveld2010diversity,
      dalvit2007traceability}.

These applications are particularly relevant to research activities
aimed at conserving local genetic resources and improving the
sustainability of livestock production systems
\cite{groeneveld2010diversity, georges2019livestock}.

High-density genotyping platforms provide detailed genomic
information, but the complete marker set may not be necessary for
every classification task.

If a substantially smaller subset of SNPs preserves most of the
breed-discriminative information contained in the original array,
the resulting reduced panel may provide a useful computational basis
for the future development of lower-density genotyping strategies.
Previous discriminant analyses have demonstrated this possibility in
both cattle and sheep, where compact marker subsets were selected for
breed or geographic assignment
\cite{dimauro2013canonical, dimauro2015sheep}.

The objective of this thesis is therefore not to claim the direct
development of a commercial diagnostic chip.
Instead, the work investigates the computational conditions under
which genomic dimensionality can be reduced while maintaining
reliable breed-classification performance.

The practical implications of the experimentally observed trade-off
between marker-panel size and predictive performance are discussed
in Chapter~\ref{chap:conclusion}.


\section{Research Questions}
\label{sec:research_questions}

Based on the computational and biological motivations described
above, the thesis addresses the following research questions:

\begin{description}

    \item[\textbf{RQ1 -- Breed classification.}]
    To what extent can machine-learning and deep
    representation-learning approaches discriminate livestock
    breeds from high-dimensional SNP genotype data?

    \item[\textbf{RQ2 -- Sample efficiency.}]
    How does breed-classification performance on unseen individuals change as the number of training individuals per breed is reduced, and what is the minimum sample size required to retain reliable generalization performance?

    \item[\textbf{RQ3 -- Marker efficiency.}]
    What is the minimum number of SNP markers required to retain
    most of the predictive performance obtained using the complete
    genomic panel?

    \item[\textbf{RQ4 -- Marker-selection strategy.}]
    Can SNP rankings derived from learned genomic representations
    and XAI methods identify informative reduced panels when
    compared with classical population-genetics and
    machine-learning baselines?

\end{description}

The primary experiments are performed on ovine genomic data.
The same experimental methodology is subsequently evaluated on an
independent caprine dataset in order to investigate whether the
observed behavior generalizes across livestock species.

Importantly, cross-species validation concerns the
\emph{methodological behavior} of the pipeline rather than the
assumption that identical SNP markers must be informative in both
species.

The relationship between the models, feature-selection procedures,
and experimental questions is summarized in
Figure~\ref{fig:thesis_overview}.
The figure provides a high-level view of the research design, while
the individual components are formalized in the following chapters.

\begin{figure}[htbp]
\centering
\begin{tikzpicture}[
    >=Latex,
    font=\small,
    box/.style={draw, semithick, rounded corners=2pt, align=center,
                minimum height=10mm, minimum width=39mm, inner sep=4pt, fill=black!2},
    key/.style={box, fill=black!7},
    arr/.style={->, semithick}
]
\node[key] (input) at (0,4.4) {High-density SNP\\genotype matrix};
\node[box] (gen) at (-3.3,2.9) {Variational\\model};
\node[box] (con) at (3.3,2.9) {Contrastive\\model};
\node[key] (repr) at (0,1.45) {Learned genomic\\representation};
\node[box] (cls) at (-3.3,0.0) {Breed\\classification};
\node[box] (xai) at (3.3,0.0) {Feature\\attribution};
\node[key] (rank) at (3.3,-1.45) {Global SNP\\ranking};
\node[key] (panel) at (3.3,-2.9) {Reduced SNP\\panels};
\node[box] (sample) at (-3.3,-2.9) {Sample-size\\analysis};
\node[box] (marker) at (3.3,-4.35) {Marker-count\\analysis};
\node[key] (cross) at (0,-5.8) {Cross-species validation\\caprine $\rightarrow$ ovine};

\draw[arr] (input) -- (gen);
\draw[arr] (input) -- (con);
\draw[arr] (gen) -- (repr);
\draw[arr] (con) -- (repr);
\draw[arr] (repr) -- (cls);
\draw[arr] (repr) -- (xai);
\draw[arr] (xai) -- (rank);
\draw[arr] (rank) -- (panel);
\draw[arr] (cls) -- (sample);
\draw[arr] (panel) -- (marker);
\draw[arr] (sample) -- (cross);
\draw[arr] (marker) -- (cross);
\end{tikzpicture}

\caption{High-level computational organization of the thesis. The workflow links representation learning to breed classification and SNP attribution, then evaluates the effects of training-set size and marker-panel size before cross-species validation. Author's elaboration.}
\label{fig:thesis_overview}
\end{figure}
